\section{Feasible sets beyond boxes}\label{sec:con}

The proof of \cref{thm:sp:main} used four properties of a mixed box.
\begin{enumerate}[label=(P\arabic*)]
\item\label{it:con:fixed} For every bag, the set of feasible values of the
coordinates outside the bag does not depend on the bag's values. The
conditional value $V_\beta$ of \eqref{eq:sp:V} is then a minimum over a fixed
set, independent of the bag's own noise, with the curvature of the objective.
\item\label{it:con:round} Every feasible point of a candidate cell can be
rounded, with its mean preserved and a second-order error, to feasible
corners that lie in the same candidate cells of every bag.
\item\label{it:con:close} Sound first-order tests identify the active
constraints near an optimizer, after which a certified strongly convex patch
contains every optimizer.
\item\label{it:con:law} The point-growth and margin events have finite-law
tails, and an exact fallback exists on the same domain.
\end{enumerate}
Arbitrary sparse affine constraints do not have these properties, and the
method cannot be extended to them without further structure:
\cref{thm:lim:constraints} shows that bounded-width affine equalities with
objective noise can encode \textsc{Subset Sum}, and \cref{ex:lim:coupled}
shows that even a single totally unimodular inequality destroys the
coordinatewise count of \cref{lem:sp:compare}. This section proves the
method for four structured classes. Each supplies these four proof steps or
replaces one of them by a proved substitute. Graph constraints and
affine-state recurrences satisfy \ref{it:con:fixed} after a reduction to a
box of retained or free coordinates, and products of simplices satisfy it
because every bag contains whole blocks. Order constraints do not satisfy
\ref{it:con:fixed}: the feasible values outside a bag depend on the bag's
values. There, a transport of feasible points supplies the semiconcave upper
bound on the conditional value that \ref{it:con:fixed} gives for boxes
(\cref{lem:con:transport}). \Cref{tab:con:summary} lists the classes.

\begin{table}[t]
\centering
\caption{Feasible-domain extensions. All use the ambient model with law
$U_{\sigma,M}$ on every original coordinate, $\log_2M\le\poly(I)$, exact
output on every draw, and the exact output of \cref{thm:count:fallback} on
rare draws. The
expected work is $C_0^{p'}\Theta^{p'}\poly(I)$, where $p'$ is the bag size
used by the algorithm; for order constraints it is the bound
\eqref{eq:con:orderbound}.}
\label{tab:con:summary}
\footnotesize
\begin{tabularx}{\textwidth}{@{}l>{\raggedright\arraybackslash}X>{\raggedright\arraybackslash}p{0.2\textwidth}l@{}}
\toprule
Result & Feasible set and structural premise & Curvature bound & Factor $\Theta$\\
\midrule
\cref{thm:con:graph}(E) & graph of an explicit polynomial map of bounded
depth over a mixed box; $p'\le p\max\{1,k^{D_{\rm g}}\}$ & retained, after
substitution & $4+(1+n/2)Lw_{\max}/(2\sigma)$\\
\cref{thm:con:graph}(I) & graph of monotone scalar equations with global
brackets; $p'\le p\max\{1,k\}$ & retained, after substitution &
$4+(1+n)Lw_{\max}/(2\sigma)$\\
\cref{thm:con:actuator} & affine-state recurrence with polynomial
actuators, any horizon; $p'=p$ & reduced, uniform in state noise &
$4+(1+n/2)Lw_{\max}/(2\sigma)$\\
\cref{thm:con:simplex} & disjoint unit simplices and integer intervals,
whole-block bags; $p'=p$ & block Hessian & $4+(2+n/2)L\max\{1,w_{\max}\}/(2\sigma)$\\
\cref{thm:con:order} & order constraints on continuous and binary
coordinates & full Hessian & $2+3n_c\bar L/(4\sigma)$, power $p_c$\\
\bottomrule
\end{tabularx}
\end{table}

Two reductions below are conditional applications of \cref{thm:sp:main} to
a family of objectives indexed by part of the noise. The constants of its
schedule must then be chosen uniformly over the family, before any noise is
drawn; choosing them after observing part of the noise would make the law
depend on the draw. The following lemma records what this requires.

\begin{lemma}[Uniform schedules]\label{lem:con:uniform}
Let $X$ be a mixed box with $n\ge1$ coordinates and a tree decomposition of
bag size $p$, and let $\mathcal H\subseteq\R^m$ be a box. For every rational
$\eta\in\mathcal H$ let $G_\eta$ be an explicit rational polynomial of degree
at most $d'$ on $X$ whose monomial support lies in a fixed set computed from
the base instance and whose factor scopes lie in the bags. Suppose that there
are polynomials $P_1,P_2$, whose degrees depend only on fixed format
constants such as degrees and depths, such that for every $b\ge1$ and every
rational $\eta\in\mathcal H$ with entries of bit length at most $b$, the
coefficients of $G_\eta$ have bit length at most $P_1(I+b)$ and $G_\eta$ is
computed from the base instance and $\eta$ in at most $P_2(I+b)$ bit
operations. Suppose also that
\begin{enumerate}[label=(\alph*)]
\item rationals $L>0$, $\kappa_2$ and $\kappa_3$, computed from the base
instance, satisfy \eqref{eq:sp:curv} and \eqref{eq:sp:kappa} for every
$G_\eta$, $\eta\in\mathcal H$, and
\item a deterministic algorithm solves every problem
$\min_X(G_\eta+\gamma^{\mathsf T}x)$ with rational $\eta\in\mathcal H$ and
$\gamma\in\Q^n$ exactly, and refines its output to precision $2^{-q}$,
within $B\poly(I+b+q)$ bit operations when the entries of $\eta$ and
$\gamma$ have bit length at most $b$; here $B\ge2$ depends only on the base
instance and $\log B\le\poly(I)$.
\end{enumerate}
Compute the schedule of \cref{def:sp:schedule} from $L$, $\kappa_2$,
$\kappa_3$, $W$, this $B$, and the constants $C_{\rm tail}$ and
$K=\max\{1,n_c3^{n_c}N_Z\max\{1,d'-1\}^{n_c}\}$ of \cref{lem:sp:tails} for
degree $d'$, and let $b$ be the largest bit length of an atom of the
resulting law $U_{\sigma,M}$. Let $\eta$ be a random rational point of
$\mathcal H$ whose entries have bit length at most $b$, with an arbitrary
law, and let $\gamma$ be independent of $\eta$ with law
$U_{\sigma,M}^{\otimes n}$. Then the algorithm of \cref{def:sp:algorithm},
with the fallback~(b), applied to $G_\eta+\gamma^{\mathsf T}x$ is correct
on every draw, runs its fallback with probability at most $1/(2B)$, and has
expected work at most $C_0^p[4+(1+n/2)Lw_{\max}/(2\sigma)]^p\poly(I+b)$,
where the degree of the polynomial factor depends on $d'$, $P_1$ and $P_2$.
Since $b\le\poly(I)$, this is $C_0^p[4+(1+n/2)Lw_{\max}/(2\sigma)]^p\poly(I)$.
\end{lemma}

\begin{proof}
Every statement of \cref{sec:sp} about one objective uses only the validity
of $L$, $\kappa_2$ and $\kappa_3$ for that objective, the tail constants,
the fallback, and the bit length of the objective's coefficients, which
enters the cost of each arithmetic operation polynomially. The constant
$C_{\rm tail}$ of \cref{lem:count:finite-tails} depends on the number of
variables, the degree and the number of atoms of the domain formula, and
$K$ on $n_c$, $N_Z$ and $d'$; neither depends on the coefficients of
$G_\eta$. All schedule quantities are therefore fixed before $\eta$ and
$\gamma$ are drawn, and $b=O(\log_2M+\operatorname{size}\sigma)$ is
polynomial in $I$. Conditionally on $\eta$, the coefficients $\gamma$ are
independent with law $U_{\sigma,M}$, so \cref{lem:sp:tails},
\cref{prop:sp:count}, \cref{lem:sp:rare} and the proof of \cref{thm:sp:main}
apply to $G_\eta$ with the same schedule, after the cost $P_2(I+b)$ of
computing $G_\eta$ and with factor values of bit length polynomial in
$P_1(I+b)$. Averaging over $\eta$ gives the claim.
\end{proof}

In the applications below, the fallback~(b) is \cref{thm:count:fallback}
applied to the \emph{original} constrained problem: an explicit polynomial
on a compact semialgebraic domain (\cref{def:count:semialg}) tilted by the
ambient noise on all coordinates. Its minimizers correspond to those of the
reduced problem, and its factor $B$ is fixed by the base instance, while the
sampled coefficients enter only through $b$.

\subsection{Global graph parameterizations}\label{sec:con:graph}

Some coupled equality constraints can be eliminated by a global chart that
keeps a set of \emph{retained} coordinates ranging over a mixed box. The
retained coordinates remain ambient coordinates with their own independent
noise. Conditioning on the noise of the eliminated coordinates then leaves
exactly the setting of \cref{sec:sp}.

\begin{definition}[Graph instance]\label{def:con:graph}
A graph instance consists of rational data as follows.
\begin{enumerate}[label=(\roman*)]
\item Retained coordinates $x$ in a mixed box $X\subseteq\R^n$ as in
\cref{def:sp:instance}\,(i), with $n\ge1$, and dependent coordinates
$y\in\R^m$ given by a chart $y=\psi(x)$ of one of two kinds.
\begin{enumerate}[label=(\Alph*)]
\item[(E)] \emph{Explicit}: $y_j=P_j(x,y_1,\dots,y_{j-1})$, $j=1,\dots,m$,
where each $P_j$ is an explicit rational polynomial of degree at most $e$
in at most $k$ of the variables (its \emph{parents}), and every chain of
dependencies has length at most $D_{\rm g}$; $e$ and $D_{\rm g}$ are fixed.
\item[(I)] \emph{Implicit}: $q_j(x_{S_j},y_j)=0$ and $y_j\in V_j=[\ell'_j,u'_j]$,
$j=1,\dots,m$, where $q_j$ is an explicit rational polynomial of degree at
most $e$, $\abs{S_j}\le k$, and rationals $\mu_j>0$ are given with
\begin{equation}\label{eq:con:bracket}
 \partial_yq_j(x,y)\ge\mu_j,\qquad q_j(x,\ell'_j)\le0\le q_j(x,u'_j)
 \qquad(x\in\widehat X_{S_j},\ y\in V_j).
\end{equation}
\end{enumerate}
The feasible set is $\mathfrak X=\{(x,\psi(x)):x\in X\}$; there are no other
constraints. Bounds on dependent coordinates in case (E) must be implied by
the chart, and in case (I) they only select the branch.
\item An objective $F_0(x,y)$ that is a sum of explicit rational polynomial
factors of degree at most $d$.
\item A tree decomposition of bag size $p$ of the ambient variables whose
bags contain every objective scope and every constraint scope, that is,
$\{y_j\}$ with its parents in case (E) and $S_j\cup\{y_j\}$ in case (I).
\item A noise half-width $\sigma>0$ and a rational $L>0$ such that
\begin{equation}\label{eq:con:curv}
 \partial_{ii}G_\eta(x)\le L,\qquad
 G_\eta(x)=F_0(x,\psi(x))+\eta^{\mathsf T}\psi(x),
\end{equation}
for all $x\in\widehat X$, all $i\in[n]$, and all $\eta\in[-\sigma,\sigma]^m$.
\end{enumerate}
The noise is ambient: $\gamma\in\Q^n$ and $\eta\in\Q^m$ have independent
coordinates with law $U_{\sigma,M}$, and the sampled objective is
$F_{\gamma,\eta}(x,y)=F_0(x,y)+\gamma^{\mathsf T}x+\eta^{\mathsf T}y$.
\end{definition}

In case (I), \eqref{eq:con:bracket} gives every $x\in\widehat X$ exactly one
root $\psi_j(x_{S_j})\in V_j$ (\cref{lem:con:chart}); the brackets and the
derivative floor must hold on the whole real hull, including between
integer labels, because the algorithm evaluates the chart at such points.
Brackets at individual query points would not establish uniqueness,
smoothness, or the validity of rounding between grid points. A dependent
coordinate whose chart has no retained argument is a constant. It is
substituted only when it is rational, that is, given by a constant
polynomial in case (E) or by a one-point interval $V_j$ in case (I). Other
constant roots, such as the root of $y^2=2$ in $[1,2]$, stay in the chart
and in the original constrained problem used by the fallback, so that no
algebraic constant enters the rational data. If preprocessing fixes every
retained coordinate, the feasible set is one point; the algorithm returns
it directly, in case (I) as the fixed retained values together with the
chart, which determines the dependent roots.

\paragraph{Status of the premises.} In case (E), $\psi$ is an explicit
polynomial map obtained by substitution, and \eqref{eq:con:curv} is
verified by monomial bounds that hold uniformly for $\abs{\eta_j}\le\sigma$,
or promised if sharper. In case (I), \eqref{eq:con:bracket} and
\eqref{eq:con:curv} are promised, or certified with their verification cost
added to the work; arbitrary polynomial positivity is not tested.
Curvature is measured after substitution. An ambient diagonal bound does
not suffice (\cref{ex:con:graph-curv}), and the dependent noise $\eta$ can
contribute to it.

\begin{lemma}[Expanded decomposition]\label{lem:con:expand}
For a dependent coordinate $y_j$ let $\Lambda(y_j)$ be the set of retained
coordinates on which it depends: its retained ancestors in case (E) and
$S_j$ in case (I); let $\Lambda(x_i)=\{x_i\}$. Replacing every bag $\beta$ by
$\beta'=\bigcup_{v\in\beta}\Lambda(v)$ gives a tree decomposition, on the same
tree, of the retained coordinates. Its bags contain the scope of every
reduced factor $f_a(x,\psi(x))$ and of every term $\eta_j\psi_j(x)$, and its
bag size is $p'\le p\max\{1,k^{D_{\rm g}}\}$ in case (E) and
$p'\le p\max\{1,k\}$ in case (I).
\end{lemma}

The proof is in \cref{app:con:graph}. The equality scopes must lie in the
supplied bags: a decomposition of the objective alone does not give running
intersection after substitution.

\begin{theorem}[Sparse optimization over global graphs]\label{thm:con:graph}
Fix $d$, $e$ and, in case (E), $D_{\rm g}$. For graph instances there is an
algorithm and a power of two $M$, computed from the base instance with
$\log_2M\le\poly(I)$, such that for the ambient model with law $U_{\sigma,M}$
on all $n+m$ coordinates the following hold.
\begin{enumerate}[label=(\roman*)]
\item On every draw the algorithm returns an exact global optimizer of
$F_{\gamma,\eta}$ on $\mathfrak X$. On ordinary draws the output is a
\emph{charted patch output} (\cref{def:con:charted}) of length $\poly(I)$
with its global proof record. On the remaining draws, of probability at most
$1/(2B)$, it is the exact output of \cref{thm:count:fallback} for
$F_{\gamma,\eta}$ on the compact semialgebraic set $\mathfrak X$, computed on
the same draw.
\item The expected bit work is at most
\[
 \text{(E)}\ \ C_0^{p'}\Bigl[4+\frac{(1+n/2)Lw_{\max}}{2\sigma}\Bigr]^{p'}\poly(I),
 \qquad
 \text{(I)}\ \ C_0^{p'}\Bigl[4+\frac{(1+n)Lw_{\max}}{2\sigma}\Bigr]^{p'}\poly(I),
\]
with $p'$ from \cref{lem:con:expand}; the global proof record obeys the same
bounds in expectation.
\item A $q$-bit evaluation returns a rational retained point $\tilde x\in X$
within distance $2^{-q}$ of the retained part $\hat x$ of the optimizer and
a rational interval $[a_q,u_q]$ of width at most $2^{-q}$ with
$a_q\le f^*\le F_{\gamma,\eta}(\tilde x,\psi(\tilde x))\le u_q$, where
$f^*=\min_{\mathfrak X}F_{\gamma,\eta}$.
In case (E) the lift $(\tilde x,\psi(\tilde x))$ is rational, satisfies every
graph equation, and lies within distance $2^{-q}$ of the optimizer. In case
(I) the lift $(\tilde x,\psi(\tilde x))$ is an exactly feasible point given
by rational retained coordinates and isolating brackets of its dependent
roots, and a rational point within distance $2^{-q}$ of the optimizer is
also returned; this rational point need not satisfy the graph equations. The
cost is $\poly(I+q)$ on ordinary draws and $B\poly(I+q)$ on fallback draws,
and its expectation is $\poly(I+q)$.
\end{enumerate}
No growth, uniqueness or strict-complementarity premise is used.
\end{theorem}

The exponents of the polynomials depend on the fixed degrees. The bound is
fixed-width polynomial in the expanded bag size $p'$; it is not a
fixed-parameter bound. Case (I) loses the factor $1+n/2$ to $1+n$ because its
costs are algebraic and are computed with certified error
(\cref{lem:con:approx}). The two cases are not comparable: (E) allows chains of
dependent coordinates of bounded depth and gives rational feasible points,
whereas (I) allows charts that are not polynomial but only one layer of
dependent coordinates. Even one cubic equation can force irrational dependent
coordinates at rational retained points, so no rational feasible point
need exist in case (I).

\begin{definition}[Charted patch output]\label{def:con:charted}
For a graph instance and a draw $(\gamma,\eta)$, a charted patch output
(\cref{def:model:outputs}\ref{it:model:charted}) is a
patch output $(S,\bar x_S,P,g_0)$ in the sense of \cref{sec:sp:setting} for
the reduced objective $G_\eta+\gamma^{\mathsf T}x$ on $X$, verified with
\eqref{eq:sp:patch}, together with the chart: the polynomials $P_j$ in case
(E); the equations $q_j$, the intervals $V_j$ and the bounds $\mu_j$ in case
(I). It denotes the point $(\hat x,\psi(\hat x))$, where $\hat x=(\bar x_S,
\hat y)$ and $\hat y$ is the unique minimizer of $G_\eta(\bar x_S,\cdot)+
\gamma^{\mathsf T}(\bar x_S,\cdot)$ on $P$. Equivalently, $\hat x$ is the
unique solution of the box KKT system of this reduced problem, and in case
(I) the dependent coordinates are the unique roots of $q_j(\hat x_{S_j},y_j)=0$
in $V_j$, with graph multipliers
$\lambda_j=-(\partial_{y_j}F_0+\eta_j)/\partial_yq_j$ at that point.
\end{definition}

The descriptor determines one point because the reduced objective is
strongly convex on $P$, and the global proof record shows that this point is
the global optimizer. No positive definiteness of the ambient Hessian is
claimed or needed.

\paragraph{The explicit case.}
By \cref{lem:con:expand}, the reduced objective $G_\eta+\gamma^{\mathsf T}x$
is an explicit polynomial of degree at most $de^{D_{\rm g}}$ on $X$ whose
factor scopes lie in the bags $\beta'$. Substitution of a dependent-noise
vector $\eta$ with entries of bit length $b$ produces coefficients of bit
length polynomial in $I+b$, in time polynomial in $I+b$, with exponents that
depend only on $d$, $e$ and $D_{\rm g}$; repeated products of supplied
coefficients can make these lengths exceed $I+b$, but only by a fixed
polynomial. The pair $(x,\psi(x))\mapsto x$ is a bijection from
$\mathfrak X$ onto $X$ that preserves objective values, so the sampled problem
is equivalent to minimizing $G_\eta+\gamma^{\mathsf T}x$ on $X$. Conditionally
on $\eta$, the coefficients $\gamma$ are independent with law $U_{\sigma,M}$.
\Cref{lem:con:uniform} therefore applies, with $L$ from
\eqref{eq:con:curv}, $\kappa_2,\kappa_3$ from monomial bounds uniform in
$\abs{\eta_j}\le\sigma$, degree $d'=de^{D_{\rm g}}$, and the fallback of
\cref{thm:count:fallback} on $\mathfrak X$, whose domain formula consists of the
mixed box and the polynomial equations $y_j=P_j$. The rest of the proof is in
\cref{app:con:graph}.

\paragraph{The implicit case.}
Here the reduced objective is algebraic but not polynomial, and three
ingredients replace exact rational arithmetic: a chart oracle
(\cref{lem:con:chart}), a dynamic program on certified lower costs
(\cref{lem:con:approx}), and closure tests with certified derivative errors.
The tails use the original polynomial graph, and so does the fallback.

\begin{lemma}[Chart]\label{lem:con:chart}
In case (I), for every $x\in\widehat X$ each equation has exactly one root
$\psi_j(x_{S_j})\in V_j$, and $\psi_j$ is infinitely differentiable on
$\widehat X_{S_j}$, with local smooth extensions at its boundary. If
$A_j\ge1$ bounds all partial derivatives of $q_j$ of orders one to three on
$\widehat X_{S_j}\times V_j$, then every partial derivative of $\psi_j$ of
order $\nu\in\{1,2,3\}$ is bounded on $\widehat X_{S_j}$ by
\[
 \Pi_{1,j}=\frac{A_j}{\mu_j},\quad
 \Pi_{2,j}=\frac{A_j(1+\Pi_{1,j})^2}{\mu_j},\quad
 \Pi_{3,j}=\frac{A_j[(1+\Pi_{1,j})^3+3(1+\Pi_{1,j})\Pi_{2,j}]}{\mu_j}.
\]
At a rational $x\in\widehat X$, rational brackets of width $\epsilon$ for all
roots, and rational enclosures of width $\epsilon$ of $G_\eta(x)$,
$\nabla G_\eta(x)$ and $\nabla^2G_\eta(x)$, are computed in time polynomial in
$I$, the length of $x$ and $\eta$, and $\log(1/\epsilon)$.
\end{lemma}

The derivative bounds have polynomial encoding length; through the chain
rule they give rationals $\kappa_2,\kappa_3$ satisfying \eqref{eq:sp:kappa}
for every $G_\eta$, $\abs{\eta_j}\le\sigma$, and the chart Lipschitz bound
\begin{equation}\label{eq:con:Kpsi}
 \norm{(x,\psi(x))-(x',\psi(x'))}_2\le K_\psi\norm{x-x'}_2,
 \qquad K_\psi=1+\sum_j\abs{S_j}\Pi_{1,j}
\end{equation}
on $\widehat X$. The proof of \cref{lem:con:chart} is in \cref{app:con:graph}.

\begin{lemma}[Certified lower costs]\label{lem:con:approx}
In case (I), at level $j$ compute for every row of every bag a rational lower
cost $\underline\phi_\beta$ with
$\underline\phi_\beta\le\phi_\beta\le\underline\phi_\beta+\delta_j$, where
$\phi_\beta$ is the sum of the reduced terms assigned to $\beta$ and
$\delta_j=E_j/N_{\mathfrak B}$ with $N_{\mathfrak B}$ the number of bags. Run
the dynamic program on the lower costs, update the incumbent with
$U_j=\min\{U_{j-1},\underline m_j+E_j\}$, where $\underline m_j$ is the
lower-cost minimum, initialize $U_{-1}=+\infty$, and store the feasible
implicit lift associated with the chosen upper bound. Update this bound
and point before pruning. Retain a cell when $\underline q_C-E_j\le U_j$, with
$\underline q_C$ the least lower-cost min-marginal over its corners. Then on
every draw all global optimizers are retained, every point removed at level
$j$ has value above $U_j\ge f^*$, and every retained cell has a corner
$v=y_\beta$ of an allowed grid point $y$ with
$G_\eta(y)+\gamma^{\mathsf T}y\le f^*+4E_j$.
\end{lemma}

The proof (\cref{app:con:graph}) is that of \cref{prop:sp:prune} with the
error $D_j=N_{\mathfrak B}\delta_j\le E_j$ of the summed lower costs. The
witness tolerance $4E_j$ replaces $2E_j$; \cref{lem:sp:compare} then gives
intervals of length $La_i+8E_j/a_i\le(1+n)La_i$, which explains the factor in
case (I). Exact comparisons of algebraic values are never made: the dynamic
program is exact rational arithmetic on certified rational costs, and
different bags may use different brackets of the same root without affecting
global consistency.

The closure test of \cref{def:sp:closure} is used with certified
derivatives. In step~(3), a rational $\hat g$ with
$\norm{\hat g-\nabla(G_\eta+\gamma^{\mathsf T}x)(c)}_\infty\le\epsilon_g$
replaces the exact gradient, and a continuous coordinate is recorded at its
lower bound when $\hat g_i-\kappa_2r-\epsilon_g>0$ and at its upper bound
when $\hat g_i+\kappa_2r+\epsilon_g<0$. In step~(5), a rational symmetric
$\hat H$ with
$\norm{\hat H-\nabla^2_{\bar S\bar S}G_\eta(\bar x_S,c^P)}_2\le\epsilon_H$
replaces the exact Hessian; it is obtained from enclosures of width at most
$\epsilon_H/n$ of the entries by taking midpoints and symmetrizing, since a
symmetric error matrix with entries at most $\epsilon_H/n$ in absolute value
has spectral norm at most $\epsilon_H$. The test is
$\hat H-(\kappa_3r^P+\epsilon_H+g_0)I\succ0$. With
$\epsilon_g=\tau/8$ and $\epsilon_H=g_0/8$, these tests are sound on every
draw and succeed under the hypotheses of \cref{prop:sp:closure-stop}
(\cref{app:con:graph}).

\begin{example}[Two necessary premises]\label{ex:con:graph-curv}
\emph{Curvature after substitution.} Under the equality $y=x$, the ambient
polynomial $F_0(x,y)=Hxy+\varepsilon(x^2+y^2)$ has both ambient diagonal
second derivatives equal to $2\varepsilon$, while its pullback
$(H+2\varepsilon)x^2$ has second derivative $2H+4\varepsilon$. The chart is
linear and well conditioned, yet the ambient cross derivative becomes
retained curvature. \emph{Retained ambient coordinates.} Under the
parameterization $(x_1,x_2)=(t_1+t_2,t_1+2t_2)$, $t\in[0,1]^2$, which has
determinant one and a bounded inverse, independent ambient coefficients pull
back to $c_1=\gamma_1+\gamma_2$ and $c_2=\gamma_1+2\gamma_2$. Under the grid
law, the event $c_2=3\sigma$ forces $\gamma_1=\gamma_2=\sigma$, so $c_1=2\sigma$
is conditionally deterministic: a singleton receives conditional probability
one instead of at most $1/M$, and the conditional count of
\cref{lem:sp:compare} fails. This does not show that the parallelogram is
hard; it shows that a bounded inverse alone does not preserve the noise model.
\end{example}

\begin{example}[A connected chain of products]\label{ex:con:chain}
Let $x\in[0,1]^n$, where some coordinates may be binary, and impose
$y_i=x_ix_{i+1}$ for $i<n$, with
$F_0=\sum_{i<n}(y_i-a_i)^2+\lambda\sum_ix_i^2$, $\lambda\ge0$. The constraints
form a connected chain of triangles $\{x_i,x_{i+1},y_i\}$, not a product of
separate blocks. Substitution gives a quartic with path bags
$\{x_i,x_{i+1}\}$, and $\partial_{ii}G_\eta=2x_{i-1}^2+2x_{i+1}^2+2\lambda\le
4+2\lambda$ for all targets $a_i$ and all $\eta$, since the dependent noise
terms $\eta_ix_ix_{i+1}$ have zero pure second derivatives. Thus
$L=4+2\lambda$, independently of $n$, and case (E) applies with $p'=2$.
\end{example}

\subsection{Affine-state recurrences with polynomial actuators}\label{sec:con:actuator}

Bounded-depth substitution does not cover dynamics over a long horizon,
because the states depend on all earlier controls. For affine state
dependence and affine state costs, an adjoint recursion eliminates the
states while adding only unary terms.

\begin{definition}[Actuator instance]\label{def:con:actuator}
For a horizon $N\ge1$ and $k=0,\dots,N-1$, the states and controls satisfy
\begin{equation}\label{eq:con:dyn}
 s_{k+1}=a_ks_k+\chi_k(u_k),\qquad\abs{a_k}\le a<1,
\end{equation}
with rational $a_k$ and $a$, explicit rational univariate polynomials $\chi_k$
of degree at most $d$, nonempty closed rational state intervals $I_k$,
$k=0,\dots,N$, and control domains $U_k$ that are closed rational intervals
or native integer intervals. Box invariance holds:
\begin{equation}\label{eq:con:invariant}
 a_kI_k+\chi_k(\operatorname{hull}U_k)\subseteq I_{k+1}\qquad(k<N).
\end{equation}
There are no other constraints. The free variables are
$z=(s_0,u_0,\dots,u_{N-1})$; a fixed initial state is allowed. The objective
is $F_0=q(z)+\sum_{k=1}^Nc_ks_k$, where $q$ is a sum of explicit rational
polynomial factors of degree at most $d$ whose scopes lie in the bags of a
supplied tree decomposition of the free variables with bag size $p$. The
noise is ambient on all original coordinates, states and controls.
\end{definition}

Condition \eqref{eq:con:invariant} is verified in polynomial time at fixed
degree, by comparing the rational interval endpoints with the extreme values
of $\chi_k$ at the endpoints of $\operatorname{hull}U_k$ and at the real roots
of $\chi_k'$, using univariate root isolation. It makes the state intervals
redundant: every free point has exactly one feasible trajectory. If
preprocessing fixes every free variable, that is, the initial state is
fixed and every control domain is a single point, the unique feasible
trajectory is rational and the algorithm returns it with its sampled value
directly.

Condition on the state coefficients $\eta=(\gamma_{s_1},\dots,\gamma_{s_N})$
and define the adjoints $\lambda_N=c_N+\eta_N$ and
$\lambda_k=c_k+\eta_k+a_k\lambda_{k+1}$ for $k=N-1,\dots,1$. Substituting
\eqref{eq:con:dyn} backward from $k=N$ gives the identity
\begin{equation}\label{eq:con:adjoint}
 \sum_{k=1}^N(c_k+\eta_k)s_k=a_0\lambda_1s_0+\sum_{k=0}^{N-1}\lambda_{k+1}\chi_k(u_k)
\end{equation}
on every trajectory. The sampled problem is therefore equivalent to
minimizing
\[
 Q_\eta(z)+\xi^{\mathsf T}z,\qquad
 Q_\eta(z)=q(z)+a_0\lambda_1s_0+\sum_{k=0}^{N-1}\lambda_{k+1}\chi_k(u_k),
\]
over the free box, where $\xi=(\gamma_{s_0},\gamma_{u_0},\dots,\gamma_{u_{N-1}})$
are independent with law $U_{\sigma,M}$ conditionally on $\eta$. The added
terms are unary, so $Q_\eta$ has the same degree and the same tree
decomposition as $q$, for every horizon. With $C_s\ge\max_k\abs{c_k}$,
induction gives $\abs{\lambda_k}\le\Lambda:=(C_s+\sigma)/(1-a)$ for every
conditioned draw. Let $L_q$ be a rational upper bound on the pure second
derivatives $\partial_{ii}q$ on the free hull; it may be negative. If
$G_\nu\ge\max_k\sup_{\operatorname{hull}U_k}\abs{\chi_k^{(\nu)}}$, then
\begin{equation}\label{eq:con:Lact}
 L=\max\{1,\ L_q+\Lambda G_2\}
\end{equation}
is positive and bounds the coordinate curvature of every $Q_\eta$, because
the added terms contribute at most $\abs{\lambda_{k+1}}\abs{\chi_k''}\le
\Lambda G_2$ to $\partial_{u_ku_k}Q_\eta$ and nothing to
$\partial_{s_0s_0}Q_\eta$. Analogous bounds with $\Lambda G_2$ and
$\Lambda G_3$ give uniform $\kappa_2$ and $\kappa_3$. The adjoints are
computed by $N$ affine steps, so their denominators divide products of the
denominators of the $a_k$, $c_k$ and $\eta_k$, and their bit lengths are
bounded by sums of the lengths of these numbers plus the length of
$\Lambda$; the coefficients of $Q_\eta$ therefore have
bit length polynomial in $I+b$, as \cref{lem:con:uniform} requires.

\begin{theorem}[Affine-state recurrences with polynomial actuators]\label{thm:con:actuator}
Fix $d$. For actuator instances, with $L$ as in \eqref{eq:con:Lact} or any
uniform valid bound, there is an algorithm and a power of two $M$, computed
from the base instance with $\log_2M\le\poly_d(I)$, such that for the ambient
model with law $U_{\sigma,M}$ the following hold, where $n$ is the number of
free variables after substituting fixed ones. The all-fixed branch
returns its rational trajectory directly; the displayed width bound
concerns the remaining nonempty free box.
\begin{enumerate}[label=(\roman*)]
\item On every draw the algorithm returns an exact global optimizer, as a
patch output in the free variables together with the recurrence
\eqref{eq:con:dyn}, or, with probability at most $1/(2B)$, as the exact
output of \cref{thm:count:fallback} for the sampled objective on the set of
feasible trajectories, described by \eqref{eq:con:dyn} and the variable
bounds.
\item The expected bit work is at most
$C_0^p[4+(1+n/2)Lw_{\max}/(2\sigma)]^p\poly_d(I)$.
\item A $q$-bit evaluation returns a rational free point and its exactly
feasible rational trajectory, within distance $2^{-q}$ of the optimal
trajectory, and a rational interval of width $2^{-q}$ containing the optimal
value, in expected bit work $\poly_d(I+q)$.
\end{enumerate}
\end{theorem}

The proof is in \cref{app:con:actuator}. The bound holds for every horizon,
because the reduced decomposition is the free-variable decomposition. The
restrictions are substantive. Nonlinear state dependence or nonlinear state
costs raise the degree or couple many controls under elimination; for
instance $s_{k+1}=s_k/2+u_k^2/4$ with terminal cost $s_N^2$ produces the
monomial $u_i^2u_l^2$ for every pair of controls. Additional path or terminal
constraints destroy the free product domain. The next example shows why the
equalities cannot simply be kept as constraints in the sparse algorithm.

\begin{example}[Predecessor fibers]\label{ex:con:dyn}
Let $s_1=u_0^2$ and $s_2=u_1^2$ with $s_0,s_1,s_2\in[0,1]$ and
$u_0,u_1\in[-1,1]$, linear objective $\gamma^{\mathsf T}(s,u)$, and bags
$\{s_0,u_0,s_1\}$ and $\{s_1,u_1,s_2\}$. Conditioning on the coefficients
outside the second bag, its conditional value at $(s_1,u_1,s_2)=(t,0,0)$ is
$\min\{0,\gamma_{s_0}\}+\gamma_{s_1}t-\abs{\gamma_{u_0}}\sqrt t$, because
$u_0=\pm\sqrt t$. Its second difference at $0,h,2h$ is
$\abs{\gamma_{u_0}}(2-\sqrt2)\sqrt h$, which is not $O(h^2)$. Under the grid
law $\gamma_{u_0}\ne0$ on every draw, since $M$ is even. No finite
coordinate curvature bound holds for this conditional value, so the count of
\cref{lem:sp:compare} does not apply, although the dynamics are box invariant
and smooth. The elimination of \eqref{eq:con:adjoint} avoids predecessor
fibers altogether.
\end{example}

\begin{example}[A stateful nonlinear actuator]\label{ex:con:actuator}
On $[-1,1]$, the transition $s_{k+1}=\frac14s_k+\frac14u_k+\frac18u_k^2$ maps
$[-1,1]\times[-1,1]$ into $[-\frac38,\frac58]$, so \eqref{eq:con:invariant}
holds. Continuous or integer controls in $[-1,1]$, affine state costs, and
any coupled nonconvex control cost $\sum_\beta q_\beta(s_0,u_\beta)$ with the
supplied decomposition satisfy \cref{def:con:actuator}.
\end{example}

\subsection{Products of simplices}\label{sec:con:simplex}

Resource constraints inside disjoint blocks are not eliminated by a chart: a
stick-breaking map from a box onto a simplex makes the ambient linear noise
nonlinear. Instead, feasible rounding and the count are adapted to the faces
of the simplices.

\begin{definition}[Simplex instance]\label{def:con:simplex}
The continuous coordinates are partitioned into blocks $b$; each block is an
\emph{inequality block} with feasible set $\Delta_b=\{x_b\ge0,\
\sum_{i\in b}x_i\le1\}$, or an \emph{equality block} with
$\Delta_b^==\{x_b\ge0,\ \sum_{i\in b}x_i=1\}$ and $\abs b\ge2$. The remaining
coordinates are native integer intervals as in \cref{def:sp:instance}, and
there is at least one coordinate. The
objective $F_0$ is a sum of explicit factors of degree at most $d$, and a
supplied tree decomposition contains every factor scope and \emph{whole
blocks}: a bag containing a coordinate of $b$ contains all of $b$. Its
largest number of scalar coordinates is $p$. A rational $L>0$ satisfies
\begin{equation}\label{eq:con:blockcurv}
 \nabla^2_{bb}F_0\preceq LI\ \text{for every block }b,\qquad
 \partial_{ii}F_0\le L\ \text{for every integer coordinate }i,
\end{equation}
on the continuous hull of the domain. The noise is ambient on all
coordinates. Let $w_{\max}$ be the largest integer width, or $1$ if there is
none, and $W=n_c+\sum_{i\in I_Z}w_i$.
\end{definition}

The bound \eqref{eq:con:blockcurv} on the block Hessian is stronger than a
bound on its diagonal entries, and it is needed: rounding within a block is
correlated. It is verified, for example, by the row-sum bound
$\max_{i\in b}\sum_{k\in b}\sup\abs{\partial_{ik}F_0}$, or promised. The
derivative bounds $\kappa_2,\kappa_3$ of \eqref{eq:sp:kappa} are taken on the
enclosing coordinate box $[0,1]^{n_c}\times\widehat X_Z$, not only on the
product of simplices, because the closure tests evaluate derivatives at
midpoints of coordinate hulls, which can lie outside a simplex. Whole-block
bags are essential: they make the domain outside a bag independent of the
bag's values. A model with simplex constraints in the scalar primal graph
satisfies this after enlarging bags by at most the largest block size; a
single resource constraint over all variables gives no width advantage.

\begin{theorem}[Products of simplices]\label{thm:con:simplex}
Fix $d$. For simplex instances there is an algorithm and a power of two $M$,
computed from the base instance with $\log_2M\le\poly_d(I)$, such that for
the ambient model with law $U_{\sigma,M}$:
\begin{enumerate}[label=(\roman*)]
\item on every draw the algorithm returns an exact global optimizer: on
ordinary draws a relative patch output of length $\poly_d(I)$ (fixed integer
labels, fixed coordinates, recorded tight budgets, a rational relative
polytope $P$, and a verified modulus $g_0$ with
$v^{\mathsf T}\nabla^2F_\gamma v\ge g_0\norm v_2^2$ on $P$ for all directions
$v$ of its affine hull), or a rational point, with its global proof record;
on the remaining draws, of probability at most $1/(2B)$, the exact output of
\cref{thm:count:fallback} on the same domain;
\item the expected bit work is at most
$C_0^p[4+(2+n/2)L\max\{1,w_{\max}\}/(2\sigma)]^p\poly_d(I)$;
\item a $q$-bit evaluation, with an exactly feasible rational point, costs
$\poly_d(I+q)$ on ordinary draws and has expectation $\poly_d(I+q)$.
\end{enumerate}
\end{theorem}

The proof is in \cref{app:con:simplex}. We describe its four ingredients.

\emph{Cells and rounding.} The mesh $h_j=s2^{-j}$ of \cref{sec:sp:cells} is
shared with the integer coordinates, so it can exceed one. While $h_j\ge1$,
each block has one cell, its simplex, whose corners are the vertices of the
simplex. For $h=h_j=1/m<1$ the cells of a block are the nonempty sets
$C=\Delta_b\cap\prod_{i\in b}[hk_i,h(k_i+1)]$ with $k_i\in\{0,\dots,m-1\}$,
or the same with $\Delta^=_b$, and the corners of $C$ are its points in
$(h\Z)^b$. After scaling, $C$ is the unit cube intersected with
$\sum_iz_i\le m-\sum_ik_i$ (or equality), whose right-hand side is an
integer, so every vertex of $C$ is a cube corner
(\cref{lem:con:simplex-round}). Every point of a cell is therefore a convex
combination of its corners, which gives mean-preserving feasible rounding
with $\E\norm{Y_b-x_b}^2\le\abs bh_j^2/4$; coordinates on grid hyperplanes
stay fixed, so all candidate lists are preserved. Rounding blocks
independently and using \eqref{eq:con:blockcurv} gives the allowance
$E_j=nLh_j^2/8$, and \cref{prop:sp:prune} holds verbatim.

\emph{Count by faces.} Classify a block grid point $v_b$ by its smallest face
of the simplex: its positive support $\mathsf S$ and whether its budget is
tight. On a slack face the directions $e_i$, $i\in\mathsf S$, are feasible
in both signs; on a tight face, with the least index $\iota\in\mathsf S$ as
anchor, the directions $e_i-e_\iota$, $i\in\mathsf S\setminus\{\iota\}$, are.
Conditioning also on the anchor coefficients, each direction confines one
remaining independent coefficient to an interval of length at most
$Lh\norm d^2+4E_j/h\le(2+n/2)Lh$. A face of dimension $r$ has at most $h^{-r}$
relative-interior grid points, which cancels the probability $h^r$
(\cref{lem:con:simplex-count}). The face classification, not a coordinate
classification, is what makes the count independent of $h$.

\emph{Closure on original faces.} For an inequality block, a positive partial
derivative throughout the hull forces $x_i=0$; a negative one forces the
budget to be tight; and a positive difference $\partial_iF-\partial_lF$
forces $x_i=0$. For an equality block only the last test is valid: decreasing
one coordinate alone is infeasible, and a strictly convex objective can have
positive partial derivatives at an interior minimizer of a probability
simplex. The margins concern \emph{active inequalities} only: zero
coordinates, and tight budgets of inequality blocks. The budget of an
equality block is imposed structurally; its multiplier is unrestricted in
sign and no margin is required of it. For $F_0=0$ on $x_1+x_2=1$, $x\ge0$,
the multiplier of the budget in the sign convention of
\cref{lem:con:simplex-close} is $-\min\{\gamma_1,\gamma_2\}$, which is
negative whenever both coefficients are positive, although the point growth
is positive whenever $\gamma_1\ne\gamma_2$; a margin requirement on it would
fail with probability close to $1/4$. At an optimizer whose active
inequalities have multipliers larger than $\tau$, the tests detect every
active inequality without a lower bound on positive coordinates: in an
inequality block with a tight budget, the derivative of a positive
coordinate is minus the budget multiplier, and in every block with a tight
or equality budget, the multiplier of a zero coordinate is its derivative
difference from a positive coordinate, in which the budget multiplier
cancels. The patch test is
$Z^{\mathsf T}\nabla^2F_\gamma(c^P)Z-(\kappa_3r^P+g_0)Z^{\mathsf T}Z\succ0$ for a
rational tangent basis $Z$ of the patch, a rational point $c^P$ of it, and the
largest hull width $r^P$; subtracting a multiple of $Z^{\mathsf T}Z$ rather
than of $I$ keeps the curvature in the metric of the original coordinates
(\cref{lem:con:simplex-close}).

\emph{Margins.} The face stationarity equations depend only on tangent noise,
the differences $\gamma_i-\gamma_\iota$, and the multiplier of each active
inequality depends with slope $\pm1$ on one further original coefficient.
The transformed coefficients are not independent under the grid law, and
the proof does not treat them as independent: it counts transformed tuples,
with at most $2M-1$ values per difference (\cref{lem:con:simplex-tail}).

\subsection{Order constraints with continuous and binary coordinates}\label{sec:con:order}

Order constraints $x_i\le x_l$ may couple the whole instance; they are not
confined to blocks, so \ref{it:con:fixed} fails. Instead, the conditional
value is controlled by transporting feasible points, and the active
constraints are identified by linear programming.

\begin{definition}[Order instance]\label{def:con:order}
Let $[n]=I_C\cup I_Z$, $n_c=\abs{I_C}\ge1$, $n_z=\abs{I_Z}$, and let $E$ be a
set of $m$ directed edges on $[n]$; cycles are allowed. The feasible set is
\[
 X=\{x\in[0,1]^n:\ x_i\le x_l\text{ for }(i,l)\in E,\ x_i\in\{0,1\}
 \text{ for }i\in I_Z\}.
\]
The objective $F_0$ is a sum of explicit factors of degree at most $d$, and
a supplied tree decomposition contains every factor scope and every edge.
Let $p$ be its bag size and $p_c$ the largest number of continuous
coordinates in a bag. A rational $\bar L\ge1$ satisfies
$\nabla^2F_0\preceq\bar LI$ on $[0,1]^n$. The noise is ambient on all
coordinates.
\end{definition}

The zero vector is feasible. If $n_c=0$, ordinary finite-state dynamic
programming over binary bag rows solves the instance exactly for every
$\gamma$ in $2^{O(p)}\poly_d(I)$ bit work, and no noise is needed. The full
Hessian bound $\bar L$ is needed because feasible rounding of ordered
coordinates is correlated; it is verified by a row-sum bound or promised.

\begin{theorem}[Order constraints]\label{thm:con:order}
Fix $d$. For order instances with $n_c\ge1$ there is an algorithm and a
power of two $M$, computed from the base instance with
$\log_2M\le\poly_d(I)$, such that for the ambient model with law
$U_{\sigma,M}$:
\begin{enumerate}[label=(\roman*)]
\item on every draw the algorithm returns an exact global optimizer: on
ordinary draws a binary assignment together with a relative patch output on
the continuous slice, consisting of fixed coordinates, a block structure
$x_{I_C}=Dy+\bar x$ with disjoint zero-one columns, a rational box and the
remaining order inequalities in $y$, and a modulus $g_0$ with verified
$\nabla^2_yF_\gamma(Dy+\bar x)\succeq g_0D^{\mathsf T}D$ on the patch, or a
rational point if no block remains; on the
remaining draws, of probability at most $1/(2B)$, the algebraic fallback
output;
\item the expected bit work is at most
\begin{equation}\label{eq:con:orderbound}
 C_0^p\,p_c!\,(p_c+1)\Bigl[2+\frac{3n_c\bar L}{4\sigma}\Bigr]^{p_c}\poly_d(I);
\end{equation}
\item a $q$-bit evaluation, with an exactly feasible rational point, costs
$\poly_d(I+q)$ on ordinary draws and has expectation $\poly_d(I+q)$.
\end{enumerate}
\end{theorem}

The number of binary coordinates is unrestricted, and they do not enter the
bracket. For $n_z=0$ the bound is
$C_0^p\,p!\,(p+1)[2+3n\bar L/(4\sigma)]^p\poly_d(I)$. The factorial is a
function of $p$ alone; the bound is fixed-width polynomial, not
fixed-parameter, because the bracket contains $n_c$. Only binary integer
coordinates are covered: transporting points with arbitrary integer levels
would need fixed interior knots and a different count. The proof is in
\cref{app:con:order}; its ingredients are as follows.

\emph{Rounding.} Continuous coordinates use dyadic cells of width
$h_j=2^{-j}$, and binary coordinates the cell $\{0,1\}$ at level $0$ and
singletons afterwards. Rounding all continuous coordinates by one common
threshold, $Y_i=h(\lfloor x_i/h\rfloor+\ind\{U<\operatorname{frac}(x_i/h)\})$
with $U$ uniform on $[0,1]$, is a nondecreasing map for each value of $U$.
It therefore preserves every order constraint and every binary value, fixes
grid nodes, and stays in every candidate cell. It is mean preserving with
$\E\norm{Y-x}^2\le n_ch^2/4$, and the full Hessian bound gives the allowance
$E_j=n_c\bar Lh_j^2/8$. The threshold is a proof device; the algorithm does
not sample it. A similar mean-preserving coupling of ordered coordinates by
one quantile appears in isotonic optimization \citep{bach2018-efficient-algorithms-for-non-convex}.

\emph{Transport and count.} Fix a bag, a binary bag row, and the noise
outside the continuous bag coordinates. For a feasible continuous bag tuple
$v$ in the relative interior of a face of a standard order simplex
$\{0\le v_{\pi(1)}\le\dots\le v_{\pi(c)}\le1\}$, take a conditional minimizer
and apply to all its coordinates the nondecreasing piecewise-affine map that
moves the distinct interior values of $v$ to those of a target $w$ in the
same closed face while fixing $0$ and $1$. The result is feasible, keeps every
binary value, has bag part $w$, and is affine in $w$; along the edge
directions of the face every coordinate moves at rate in $[0,1]$, so the
transported objective is an upper support for the conditional value with
directional curvature at most $n_c\bar L$ (\cref{lem:con:transport}). The
transported fiber need not equal the true fiber, and the projected mixed
domain need not be convex; an upper support is what the comparison needs.
Counting relative-interior grid points of the faces of the $p_c!$ order
simplices, as in \cref{lem:sp:compare}, gives the bracket of
\eqref{eq:con:orderbound} (\cref{prop:con:order-count}).

\emph{Closure by linear programming.} When every binary hull is a singleton,
all optimizers lie in one continuous slice, a polytope with zero-one
vertices. With $g$ the gradient at the hull midpoint and $\delta=\kappa_2r$
a bound on its error, a proposed equality $x_i=x_l$ is certified when the
least value of $g^{\mathsf T}x$ on the slice face $\{x_i=0,x_l=1\}$ exceeds
the least value on the slice by more than $n_c\delta$; bound equalities use the
faces $\{x_i=1\}$ and $\{x_i=0\}$. These are rational linear programs. At
every optimizer the true slice gradient exposes a face containing it, and if
the equality failed some zero-one vertex of that face would lie in the
opposite face, which bounds the computed gap by $n_c\delta$
(\cref{lem:con:lpgap}). The certified equalities merge coordinates into
blocks, and all other order inequalities are kept in the patch, so no lower
bound on inactive slacks is needed.

\begin{example}[Binary labels must be fixed before exposure]\label{ex:con:premature}
On $0\le x\le z\le1$ with $z$ binary, the objective
$(x-\frac12)^2+10z^2-11z$ has the unique mixed optimizer $(\frac12,1)$ with
value $-1$ and point-growth modulus one. Its full gradient there is $(0,9)$.
Over the continuous relaxation, the least value of $9z$ is $0$, while on the
face $\{x=0,z=1\}$ opposite to the equality $x=z$ it is $9$. The relaxed test
would certify $x=z$, which is false. After $z=1$ is fixed, the slice gradient
is zero and no equality is certified. The algorithm therefore applies the
linear-programming test only after all binary hulls are singletons.
\end{example}

\subsection{What is not covered}\label{sec:con:scope}

General totally unimodular constraints are not covered. For
$P=\{x\in[0,1]^n:Ax\le b\}$ with $A$ integral and totally unimodular and $b$
integral, aligned feasible rounding holds: after scaling a cell of width
$h=2^{-j}$, the cell polytope $\{z\in[0,1]^n:Az\le b/h-Ak\}$ has an integral
right-hand side and hence integral vertices
\citep{hoffman1956-integral-boundary-points-of-convex}, so every feasible cell point is a mean of feasible
corners, and a full Hessian bound $\bar L$ gives the allowance
$n\bar Lh^2/8$. With every constraint inside a bag, the pruning of
\cref{prop:sp:prune} is then sound. What is missing is a count of retained
cells whose constants are controlled by the input: \cref{ex:lim:coupled}
shows that coordinate comparisons fail for a single totally unimodular
inequality, and the arguments of this paper provide neither such a count
nor closure and tail estimates for general totally unimodular systems. Order
constraints are the totally unimodular case in which the transport count,
the linear-programming closure and the tails are all complete. Overlapping resource constraints, affine
images of simplices, general integer ranges in order systems, extra
constraints that cut a graph chart, and nonlinear state recurrences are also
outside the theorems of this section.

\paragraph{Relation to earlier work.} The geometric ingredients are
classical: order polytopes have zero-one vertices and are triangulated by
order simplices \citep{stanley1986-two-poset-polytopes}; isotonic optimization uses the
common-quantile coupling \citep{bach2018-efficient-algorithms-for-non-convex}; totally unimodular
systems have integral vertices \citep{hoffman1956-integral-boundary-points-of-convex}; and implicit
elimination of states inside validated global search for dynamic models is
established \citep{wilhelm2019-global-optimization-of-stiff-dynamical}. The new element is again the
composition: an expected near-optimal row count under one finite
independent-noise law for each domain class, sound closure tests adapted to
the domain, and an exact same-draw fallback with a polynomial-time
evaluable output.
